% ID: SEL-D0028-OA-E04
% ESTADO: cotejada. Segunda pasada visual completa y borrador compilado revisados.
% Fuente: Selectividad/matematicas_aplicadas/suplente junio/Examen-Matemáticas Aplicadas a las Ciencias Sociales.pdf
% SHA-256: a10146fce64a17e21450973c88d66d43571619781f76bc345b5a4b3dee67a16b
% Página: 1; opción A; ejercicio original 4.

A la salida de unos grandes almacenes se ha tomado una muestra aleatoria simple de 100 clientes, a los que se les ha preguntado por el gasto que han realizado, obteniéndose una media muestral de 110 euros. Se sabe que el gasto sigue una distribución Normal con desviación típica 20 euros.
\begin{enumerate}[label=\alph*)]
\item \textbf{(0.5 puntos)} \textquestiondown{}Qué distribución de probabilidad sigue la media muestral?
\item \textbf{(1 punto)} Obtenga un intervalo de confianza al 90\%, para el gasto medio de todos los clientes que han comprado ese día.
\item \textbf{(1 punto)} Si deseamos que el error máximo cometido, con el mismo nivel de confianza, sea 2 euros, \textquestiondown{}cuál ha de ser el tamaño mínimo de la muestra?
\end{enumerate}
